% =====================================================================
%  EE6423 – Bản mô tả bài toán 1 trang (M0), theo Phụ lục A của project brief
%  Biên dịch:  latexmk -xelatex EE6423_M0_MoTaBaiToan.tex
% =====================================================================
\documentclass[a4paper,11pt]{article}
\usepackage{fontspec}
\usepackage{polyglossia}
\setdefaultlanguage{vietnamese}
\setmainfont{Times New Roman}
\usepackage{amsmath}
\usepackage{unicode-math}
\setmathfont{TeX Gyre Termes Math}
\usepackage{icomma}
\usepackage[a4paper,left=2cm,right=2cm,top=1.5cm,bottom=1.5cm]{geometry}
\usepackage{enumitem}
\usepackage{tabularx,booktabs}
\usepackage{xcolor}
\pagestyle{empty}
\setlength{\parindent}{0pt}
\setlength{\parskip}{2pt}
\setlist{nosep,leftmargin=1.2em}
\newcommand{\muc}[1]{\par\smallskip\textbf{#1}\enspace}
\newcommand{\co}{\textbf{Có}}
\newcommand{\cc}{\textbf{Chưa chắc}}

\begin{document}

\begin{center}
  {\large\bfseries EE6423 – Bản mô tả bài toán (M0)}\\[2pt]
  \textbf{Thành viên:} Nguyễn Thái Sơn – 20252373M \quad·\quad Vũ Minh Đức – 20252371M
\end{center}
\vspace{-4pt}\rule{\linewidth}{0.4pt}

\muc{1. Tên bài toán:} Ước lượng tín hiệu chuyển động của người từ cảm biến hồng ngoại thụ động PYD~1378
khi có nhiễu điện lưới 50~Hz, trôi nền nhiệt và nhiễu nội tại cùng dải tần với tín hiệu.

\muc{2. Hệ và cảm biến:}
\begin{itemize}
  \item \emph{Hệ/thiết bị:} đầu báo chuyển động PIR gồm thấu kính Fresnel, cảm biến, mạch khuếch đại -- lọc analog
        (62~dB, dải thông 0,5 -- 8,5~Hz, nguồn 3,3~V, đầu ra 0 -- 3~V, mức nghỉ 1,5~V), ADC 12~bit và xử lý số.
  \item \emph{Cảm biến:} PYD~1378, đầu thu nhiệt điện hai phần tử A, B nối tiếp ngược cực, có FET lặp nguồn.
        Nguyên lý: dòng nhiệt điện $i_p = pA\,\mathrm{d}T/\mathrm{d}t$, chỉ đáp ứng khi bức xạ thay đổi.
        Độ nhạy 4~kV/W tại 1~Hz, nhiễu 20~µV$_\text{pp}$; đơn vị số đo: V.
\end{itemize}

\muc{3. Đại lượng cần ước lượng:} dạng sóng chuyển động $x[n]$ (V) do người đi qua vùng nhìn, gần tuần hoàn
0,15 -- 6,7~Hz, biên độ 0,05 -- 1,5~V tại đầu ra. Từ $x[n]$ suy ra có/không có người và hướng đi.
Số đo không cho trực tiếp $x[n]$ vì tín hiệu ở chân S chỉ vài mV và bị chồng bởi các nguồn nhiễu dưới đây.

\muc{4. Các nguồn bất định dự kiến:}
\begin{itemize}
  \item Nhiễu điện lưới 50~Hz và các hài, cảm ứng lên dây nối và cực cổng trở kháng cao (thành phần tuần hoàn).
  \item Trôi nền nhiệt (nhiệt độ phòng, nắng, luồng khí nóng) và lệch mức nghỉ do op-amp, sai số linh kiện (bias/drift).
  \item Nhiễu nội tại của cảm biến (nhiễu nhiệt, nhiễu 1/f), bị định dạng bởi dải thông nên là nhiễu màu, trùng dải tín hiệu.
  \item Nhiễu trắng của op-amp, điện trở và lượng tử hóa ADC (LSB 0,81~mV).
\end{itemize}

\muc{5. Các chế độ vận hành dự kiến:} (1) nghỉ, không có người, nền ổn định; (2) người đi qua với tốc độ và khoảng cách
khác nhau, theo hai hướng; (3) điều kiện thay đổi: nền nhiệt tăng dần, thiết bị điện gần đó bật/tắt làm biên độ 50~Hz đổi.

\muc{6. Nguồn dữ liệu và ground truth:} mô phỏng bằng Python (seed cố định, $f_s = 200$~Hz). Mô hình cảm biến lấy từ mô hình
nhiệt -- điện đã hiệu chuẩn theo datasheet PYD~1378 và kiểm chứng trong LTspice (sai lệch độ nhạy dưới 10\%).
Ground truth là $x[n]$ sạch của mô phỏng; chế độ nghỉ cho tham chiếu biết trước $x[n] = 0$.

\muc{7. Tự kiểm tra theo checklist mục 4:}\\[2pt]
\small
\begin{tabularx}{\linewidth}{@{}p{3.4cm}p{1.6cm}X@{}}
\toprule
\textbf{Dòng} & \textbf{Trả lời} & \textbf{Lý do} \\ \midrule
\multicolumn{3}{@{}l}{\textit{Hồi I – Hiểu nhiễu}} \\
Nhiễu có cấu trúc & \co & Điện lưới 50~Hz (vạch phổ), trôi nền (bướu hẹp tại 0), nhiễu nội tại màu. \\
Biết $f_s$ và chuỗi đo & \co & PYD~1378 $\to$ khuếch đại -- lọc analog $\to$ ADC 12~bit, $f_s = 200$~Hz $\to$ xử lý số. \\
\multicolumn{3}{@{}l}{\textit{Hồi II – Học từ dữ liệu}} \\
Tham số cần ước lượng & \co & Biên độ 50~Hz, bias, phương sai nhiễu, tham số tự tương quan của nhiễu nội tại. \\
Dữ liệu hiệu chuẩn & \co & Chế độ nghỉ biết trước $x[n] = 0$; ground truth đầy đủ từ mô phỏng. \\
\multicolumn{3}{@{}l}{\textit{Hồi III – Hành động}} \\
Đại lượng theo thời gian & \co & $x[n]$ thay đổi theo thời gian và chìm trong nhiễu nội tại cùng dải tần. \\
Mô hình thống kê/động học & \cc & $x[n]$ là chuỗi xung gần tuần hoàn; mô hình động học phù hợp cần kiểm chứng ở M2. \\
Kịch bản thay đổi & \co & Nền nhiệt trôi, biên độ 50~Hz đổi khi bật/tắt thiết bị, tốc độ và khoảng cách người đổi. \\
\multicolumn{3}{@{}l}{\textit{Đánh giá}} \\
Có metric & \co & RMSE của $x[n]$ ước lượng so với ground truth, SNR trước/sau xử lý. \\
$\ge 2$ phương án & \cc & Dự kiến có lọc theo thống kê dừng và lọc theo mô hình trạng thái; chọn ở M3. \\
\bottomrule
\end{tabularx}

\end{document}
